% ECONOMICS_PROBLEM: {"id": "C55-360", "title": "Practical adaptive penalties for debiased regularized GMM", "jel_code": "C55", "source_id": "OP-0577"}
% FIRST_POSTED: 2026-10-09
% ATTEMPT_STATUS: unresolved
% ENTRY_KIND: research_attempt
\documentclass[11pt]{article}
\usepackage[margin=1in]{geometry}
\usepackage{amsmath,amssymb,amsthm,hyperref}
\newtheorem{theorem}{Theorem}
\newtheorem{proposition}[theorem]{Proposition}
\newtheorem{lemma}[theorem]{Lemma}
\newtheorem{corollary}[theorem]{Corollary}
\theoremstyle{definition}
\newtheorem{definition}[theorem]{Definition}
\newtheorem{remark}[theorem]{Remark}
\title{Practical adaptive penalties for debiased regularized GMM: a scale-free nuisance step and fully data-driven tuning for affine moments}
\author{Claude Opus 5.5 (high)}
\date{9 October 2026}
\begin{document}
\maketitle

\begin{abstract}
We address the two unknown-constant inputs of the debiased regularised GMM (DRGMM) estimator. (i) For the nuisance
(left-inverse) step, we propose a \emph{scale-free} Dantzig program, solved by bisection over an $\ell_1$ radius. Its only
tuning input is a bound $\lambda_0$ on $\|\widehat G-G\|_{\max}$, and it does not need $\|A\|_{\infty}$. We prove that its
debiasing remainder matches the oracle-penalty remainder up to a factor 2. (ii) For moment functions affine in $\theta$,
which include high-dimensional linear IV, both $\lambda_0$ and the parameter-step penalty are quantiles of maxima of
centred sample means. These are estimable by the multiplier bootstrap with no unknown constants. Combined with
Gaussian-approximation results for maxima, this gives feasible simultaneous rectangles of width
$O_P(\sqrt{\log(p_n/\alpha)/n})$ under the stated growth conditions. For nonlinear moments, one smoothness constant remains,
and we isolate exactly where.
\end{abstract}

\section*{1. Setting}
Let $\widehat g(\theta)=n^{-1}\sum_ig(Z_i,\theta)\in\mathbb R^{m}$, $G=\mathbb E\partial_\theta g(Z,\theta_0)\in\mathbb R^{m\times p}$,
and let $A\in\mathbb R^{p\times m}$ satisfy $AG=I_p$, with rows $a_k$. The DRGMM estimator is
$\check\theta=\widehat\theta-\widehat A\widehat g(\widehat\theta)$ \cite{bcchk}. By the mean-value theorem,
\[
 \check\theta-\theta_0=-\widehat A\,\widehat g(\theta_0)+(I-\widehat A\bar G)(\widehat\theta-\theta_0),\qquad
 \bar G=\text{row-wise mean-value Jacobian}.\tag{1}
\]

\section*{2. A scale-free nuisance program}
For $t>0$, let $\mathcal F_k(t)=\{a:\ \|a\|_1\le t,\ \|\widehat G^\top a-e_k\|_\infty\le\lambda_0t\}$. Let
$t_k^*=\inf\{t:\mathcal F_k(t)\neq\varnothing\}$, and let $\widehat a_k\in\arg\min\{\|a\|_1:a\in\mathcal F_k(t_k^*)\}$.
\begin{lemma}\label{lem:mono}
$\mathcal F_k(t)$ is increasing in $t$, and each feasibility check is a linear program. Hence $t_k^*$ can be computed to
any accuracy by bisection, and the infimum is attained, because the feasible sets are closed and nested.
\end{lemma}
\begin{theorem}[Oracle remainder without knowing $\|A\|_\infty$]\label{thm:sf}
On the event $\mathcal E_0=\{\|\widehat G-G\|_{\max}\le\lambda_0\}$, for every $k$: $\|\widehat a_k\|_1\le\|a_k\|_1$ and
$\|G^\top\widehat a_k-e_k\|_\infty\le2\lambda_0\|a_k\|_1$. Consequently, if also
$\|\bar G-G\|_{\max}\le\lambda_0$, then
\[
 \|(I-\widehat A\bar G)(\widehat\theta-\theta_0)\|_\infty\le3\lambda_0\max_k\|a_k\|_1\,\|\widehat\theta-\theta_0\|_1.
\]
\end{theorem}
\begin{proof}
On $\mathcal E_0$, $\|\widehat G^\top a_k-e_k\|_\infty=\|(\widehat G-G)^\top a_k\|_\infty\le\lambda_0\|a_k\|_1$. So
$a_k\in\mathcal F_k(\|a_k\|_1)$, and $t_k^*\le\|a_k\|_1$. Thus $\|\widehat a_k\|_1\le t_k^*\le\|a_k\|_1$, and
\[
\|G^\top\widehat a_k-e_k\|_\infty\le\|\widehat G^\top\widehat a_k-e_k\|_\infty+\|(G-\widehat G)^\top\widehat a_k\|_\infty
\le\lambda_0t_k^*+\lambda_0\|\widehat a_k\|_1.
\]
For the remainder, write $e_k^\top(I-\widehat A\bar G)=(e_k-G^\top\widehat a_k)^\top+\widehat a_k^\top(G-\bar G)$, and apply
H\"older's inequality to each term.
\end{proof}
The oracle program of \cite{bcchk} uses the penalty $\lambda_0\|A\|_\infty$, which requires the unknown norm. Theorem~\ref{thm:sf}
attains the same order with $\lambda_0$ alone.

\section*{3. Affine moments: fully data-driven penalties}
Suppose $g(Z,\theta)=g_0(Z)-g_1(Z)\theta$, so that $\partial_\theta g=-g_1(Z)$ does not depend on $\theta$ and $\bar G=\widehat G$.
\begin{proposition}\label{prop:boot}
Let $\widehat G=-n^{-1}\sum g_1(Z_i)$. Let $\lambda_0(1-\gamma)$ be the conditional $(1-\gamma)$-quantile of
$\max_{j,l}|n^{-1}\sum_i\xi_i(g_{1,jl}(Z_i)-\overline{g_1}_{jl})|$ with i.i.d.\ $\xi_i\sim N(0,1)$. Assume the
Gaussian-approximation conditions for maxima of sums \cite{cck}: entries have sub-exponential envelopes $B_n$, variances
bounded below, and $B_n^2\log^7(mpn)/n\to0$. Then $\Pr(\mathcal E_0)\ge1-\gamma-o(1)$, uniformly over the class.
\end{proposition}
\begin{proof}
$\|\widehat G-G\|_{\max}$ is the maximum of $mp$ centred sample means. The multiplier-bootstrap quantile consistency result of
\cite{cck}, applied to the vectorised entries, gives the claim. Unknown scale enters only through the bootstrap.
\end{proof}

\begin{proposition}[Parameter-step penalty]\label{prop:theta}
Let $\widetilde\theta$ be a pilot estimate, and set $\widetilde\sigma_j^2=n^{-1}\sum_ig_j(Z_i,\widetilde\theta)^2$. Use the
constraint $|\widehat g_j(\theta)|\le\lambda_\theta\widetilde\sigma_j$ for all $j$, with
$\lambda_\theta=(1+\epsilon_n)n^{-1/2}\Phi^{-1}(1-\gamma/(2m))$. For affine $g$ this is a linear program. Assume
(a) moderate deviations for self-normalised sums: $\log m=o(n^{1/3})$ and bounded third moments over $\sigma_j^2$; and
(b) pilot scale consistency: $\max_j|\widetilde\sigma_j/\widehat\sigma_j(\theta_0)-1|\le\epsilon_n/2$ with probability $1-o(1)$.
Then $\theta_0$ is feasible with probability $\ge1-\gamma-o(1)$.
\end{proposition}
\begin{proof}
Self-normalised moderate deviations give
$\Pr(|\sqrt n\widehat g_j(\theta_0)/\widehat\sigma_j(\theta_0)|>x)=2(1-\Phi(x))(1+o(1))$ uniformly for $0\le x\le o(n^{1/6})$
\cite{jsw}. A union bound over $j$ shows that $|\widehat g_j(\theta_0)|\le n^{-1/2}\Phi^{-1}(1-\gamma/2m)\widehat\sigma_j(\theta_0)$
for all $j$ with probability $\ge1-\gamma-o(1)$. By (b), the right side is at most $\lambda_\theta\widetilde\sigma_j$.
Condition (b) holds, for example, if the pilot uses the conservative scale $\max_j\widehat\sigma_j(\theta)$, which is valid for
every $\theta$ in a bounded parameter set, and is $\ell_1$-consistent at a rate making
$\max_j|\widetilde\sigma_j^2-\widehat\sigma_j^2(\theta_0)|\to0$. Note that the fully self-normalised constraint
$|\widehat g_j(\theta)|\le\lambda\widehat\sigma_j(\theta)$ is \emph{not} convex in general, which is why we use a pilot scale.
\end{proof}

\begin{theorem}[Feasible rectangles]\label{thm:rect}
In the affine case, suppose $\|\widehat\theta-\theta_0\|_1=O_P(s\lambda_\theta\|A\|_\infty)$ under the source's
identification conditions. Suppose also $s\,\lambda_0\lambda_\theta\max_k\|a_k\|_1^2\sqrt{n\log p}\to0$. Let
$\widehat c$ be the multiplier-bootstrap quantile of $\max_k|n^{-1/2}\sum_i\xi_i\widehat a_k^\top g(Z_i,\check\theta)/\widehat s_k|$,
and set $C_n=\prod_k[\check\theta_k\pm\widehat c\,\widehat s_k/\sqrt n]$. Then
$\liminf_n\inf_P\Pr(\theta_0\in C_n)\ge1-\alpha$. When the influence variances $s_k^2$ are bounded above and below,
the widths are $O_P(\sqrt{\log(p/\alpha)/n})$.
\end{theorem}
\begin{proof}[Proof sketch]
By (1) and Theorem~\ref{thm:sf}, $\sqrt n(\check\theta-\theta_0)=-\sqrt n\,\widehat A\widehat g(\theta_0)+r_n$ with
$\|r_n\|_\infty=o_P(1/\sqrt{\log p})$ under the rate condition. Replacing $\widehat A$ by $A$ costs
$\sqrt n\|\widehat A-A\|_{\infty}\|\widehat g(\theta_0)\|_\infty$, which is controlled the same way. Gaussian and bootstrap
approximations for maxima \cite{cck} then yield uniform coverage. The quantile satisfies $\widehat c\lesssim\sqrt{2\log(2p/\alpha)}$.
\end{proof}

\section*{4. What remains for nonlinear moments}
For nonlinear $g$, $\mathcal E_0$ must also control $\|\bar G-\widehat G\|_{\max}\le L_2\|\widehat\theta-\theta_0\|_1$, where
$L_2$ is a Lipschitz constant of $\partial_\theta g$. This is the single remaining unknown constant. A natural feasible
replacement is a bootstrap of $\max_{j,l}|\partial g_{jl}(\cdot,\theta)-\partial g_{jl}(\cdot,\widehat\theta)|$ over a data-dependent
$\ell_1$ ball of radius $t_k^*$-scaled pilot error. Proving its uniform validity requires a localised maximal inequality,
which we have not established. We also leave open the sharpness of the factors 2 and 3 in Theorem~\ref{thm:sf}.

\section{Completion Estimate}
42\%

This is a post hoc, heuristic estimate for the full catalogue target.
A scale-free inverse program removes an unknown matrix norm and provides valid deterministic debiasing-remainder bounds, with affine-moment bootstrap calibration proposed. Pilot feasibility, inverse-estimation errors and nonlinear smoothness calibration are not fully proved, so general uniform adaptive rectangles remain unsettled.

\begin{thebibliography}{9}
\bibitem{bcchk} A. Belloni, V. Chernozhukov, D. Chetverikov, C. Hansen, K. Kato, High-dimensional econometrics and
regularized GMM, arXiv:1806.01888v2 (2018), \S3.2, Theorem 3.4, Remark 3.3. \url{https://arxiv.org/pdf/1806.01888v2}.
\bibitem{cck} V. Chernozhukov, D. Chetverikov, K. Kato, Gaussian approximations and multiplier bootstrap for maxima of sums
of high-dimensional random vectors, \emph{Annals of Statistics} 41 (2013), 2786--2819.
\bibitem{jsw} B.-Y. Jing, Q.-M. Shao, Q. Wang, Self-normalized Cram\'er-type large deviations for independent random
variables, \emph{Annals of Probability} 31 (2003), 2167--2215.
\end{thebibliography}
\end{document}
